%%%%%%%%% pl11.tex %%%%%%
\begin{abstract}
We introduce a notion of planar
algebra, the simplest example of which is a vector
space of tensors, closed under planar
contractions. A planar algebra with suitable
positivity properties produces a finite index
subfactor of a II${}_1$ factor, and vice versa.
\end{abstract}

\baselineskip=16 pt

\noindent
{\large\bf 0. Introduction}

\medskip At first glance there is nothing planar
about a subfactor. A factor $M$ is a unital
$*$-algebra of bounded linear operators on a
Hilbert space, with trivial centre and closed in
the topology of pointwise convergence. The factor
$M$ is of type II${}_1$ if it admits a
(normalized) trace, a linear function $tr$:
$M\to\Bbb C$ with $tr(ab)=tr(ba)$ and
tr$(1)=1$. In [J1] we defined the notion of {\it
index} $[M:N]$ for II${}_1$ factors $N\subset M$.
The most surprising result of [J1] was that
$[M:N]$ is ``quantized" --- to be precise, if
$[M:N] <4$ there is an integer $n\geq 3$ with
$[M:N]=4\cos^2\pi/n$. This led to a surge of
interest in subfactors and the major theorems of
Pimsner, Popa and Ocneanu ([PP],[Po1],[O1]). These
results turn around a ``standard invariant" for
finite index subfactors, also known variously as
the ``tower of relative commutants", the
``paragroup", or the ``$\l$-lattice". In favorable
cases the standard invariant allows one to
reconstruct the subfactor, and both the paragroup
and $\l$-lattice approaches give complete
axiomatizations of the standard invariant. In this
paper we give, among other things, yet another
axiomatization which has the advantage of
revealing an underlying {\it planar} structure not
apparent in other approaches.  It also places the standard
invariant in a larger mathematical context. In
particular we give a rigorous justification for
pictorial proofs of subfactor theorems.
Non-trivial results have already been obtained
from such arguments in [BJ1]. The standard
invariant is sufficiently rich to justify several
axiomatizations --- it has led to the discovery of
invariants in knot theory ([J2]), 3-manifolds
([TV]) and combinatorics ([NJ]), and is of
considerable interest in conformal and algebraic
quantum field theory ([Wa],[FRS],[Lo]).

Let us now say exactly what we mean by a planar
algebra. The best language to use is that of
operads ([Ma]). We define the {\it planar operad},
each element of which determines a multilinear
operation on the standard invariant.

A planar $k$-tangle will consist of the unit disc
$D$ $(=D_0)$ in $\Bbb C$ together with a finite
(possibly empty) set of disjoint subdiscs
$D_1,D_2,\dots ,D_n$ in the interior of $D$. Each
disc $D_i$, $i\geq 0$, will have an even number
$2k_i\geq 0$ of marked points on its boundary
(with $k=k_0$). Inside $D$ there is also a finite
set of disjoint smoothly embedded curves called
{\it strings} which are either closed curves or
whose boundaries are marked points of the $D_i$'s.
Each marked point is the boundary point of some
string, which meets the boundary of the
corresponding disc transversally. The strings all
lie in the complement of the interiors $\stackrel{\circ}{D_i}$
of the $D_i$, $i\geq 0$. The connected components of
the complement of the strings in 
$\dsize{\stackrel{\circ}{D}\backslash \bigcup^n_{i=1}D_i}$ 
are called regions
and are shaded black and white so that regions
whose closures meet have different shadings. The
shading is part of the data of the tangle, as is
the choice, at every $D_i$, $i\geq 0$, of a white
region whose closure meets that disc. 
The case $k=0$ is exceptional - 
there are two kinds of 0-tangle, according to
whether the region near the boundary is shaded
black or white. An example of a planar 4-tangle,
where the  chosen white regions are marked with a
$*$ close to their respective discs, is given
below.

\[
\begin{picture}(0,0)%
\epsfig{file=ps/pic1.ps}%
\end{picture}%
\setlength{\unitlength}{0.00041700in}%
%
\begingroup\makeatletter\ifx\SetFigFont\undefined
% extract first six characters in \fmtname
\def\x#1#2#3#4#5#6#7\relax{\def\x{#1#2#3#4#5#6}}%
\expandafter\x\fmtname xxxxxx\relax \def\y{splain}%
\ifx\x\y   % LaTeX or SliTeX?
\gdef\SetFigFont#1#2#3{%
  \ifnum #1<17\tiny\else \ifnum #1<20\small\else
  \ifnum #1<24\normalsize\else \ifnum #1<29\large\else
  \ifnum #1<34\Large\else \ifnum #1<41\LARGE\else
     \huge\fi\fi\fi\fi\fi\fi
  \csname #3\endcsname}%
\else
\gdef\SetFigFont#1#2#3{\begingroup
  \count@#1\relax \ifnum 25<\count@\count@25\fi
  \def\x{\endgroup\@setsize\SetFigFont{#2pt}}%
  \expandafter\x
    \csname \romannumeral\the\count@ pt\expandafter\endcsname
    \csname @\romannumeral\the\count@ pt\endcsname
  \csname #3\endcsname}%
\fi
\fi\endgroup
\begin{picture}(7094,7092)(286,-6840)
\put(2166,-2580){\makebox(0,0)[lb]{\smash{\SetFigFont{12}{14.4}{rm}$D_6$}}}
\put(4746,-2980){\makebox(0,0)[lb]{\smash{\SetFigFont{12}{14.4}{rm}$*$}}}
\put(3931,-2870){\makebox(0,0)[lb]{\smash{\SetFigFont{12}{14.4}{rm}$D_1$}}}
\put(5676,-2465){\makebox(0,0)[lb]{\smash{\SetFigFont{12}{14.4}{rm}$D_2$}}}
\put(5361,-3955){\makebox(0,0)[lb]{\smash{\SetFigFont{12}{14.4}{rm}$D_3$}}}
\put(5971,-4265){\makebox(0,0)[lb]{\smash{\SetFigFont{12}{14.4}{rm}$*$}}}
\put(4516,-5225){\makebox(0,0)[lb]{\smash{\SetFigFont{12}{14.4}{rm}$D_4$}}}
\put(4596,-5820){\makebox(0,0)[lb]{\smash{\SetFigFont{12}{14.4}{rm}$*$}}}
\put(2146,-5135){\makebox(0,0)[lb]{\smash{\SetFigFont{12}{14.4}{rm}$D_5$}}}
\put(1636,-4820){\makebox(0,0)[lb]{\smash{\SetFigFont{12}{14.4}{rm}$*$}}}
\put(3811,-43){\makebox(0,0)[lb]{\smash{\SetFigFont{12}{14.4}{rm}$*$}}}
\put(5571,-2915){\makebox(0,0)[lb]{\smash{\SetFigFont{12}{14.4}{rm}$*$}}}
\put(3166,-1580){\makebox(0,0)[lb]{\smash{\SetFigFont{12}{14.4}{rm}$*$}}}
\put(3256,-1078){\makebox(0,0)[lb]{\smash{\SetFigFont{12}{14.4}{rm}$D_7$}}}
\put(2096,-3000){\makebox(0,0)[lb]{\smash{\SetFigFont{12}{14.4}{rm}$*$}}}
\end{picture}

\]
The {\it planar operad} $\Bbb P$ is the set of all
orientation-preserving diffeomorphism classes of
planar $k$ tangles, $k$ being arbitrary. The
diffeomorphisms preserve the boundary of $D$ but
may move the $D_i$'s, $i>1$.

Given a planar $k$ tangle $T$, a $k'$-tangle $S$,
and a disk $D_i$ of $T$ with $k_i=k'$ we define
the  $k$ tangle $T\circ_i S$ by isotoping $S$ so
that its boundary, together with the marked
points, coincides with that of $D_i$, and the
chosen white regions for $D_i$ (in $T$) and $S$ share a
boundary segment.  The strings may then be joined
at the boundary of $D_i$ and smoothed. The
boundary of $D_i$ is then removed to obtain the
tangle $T\circ_i S$ whose diffeomorphism class
clearly depends only on those of $T$ and $S$.
This gives $\Bbb P$ the structure of a coloured
operad, where each $D_i$ for $i>0$ is assigned
the colour $k_i$ and composition is only allowed
when the colours match. There are two distinct
colours for $k=0$ according to the shading near
the boundary. The $D_i$'s for $i\geq 1$ are to be
thought of as inputs and $D_0$ is the output. (In
the usual definition of an operad the inputs are
labelled and the symmetric group $S_n$ acts on
them. Because of the colours, $S_n$ is here replaced by
$S_{n_1}\x S_{n_2}\x\dots\x S_{n_p}$ where  $n_j$ is the
number of internal discs coloured $j$.
Axioms for such a coloured operad could be given
along the lines of [Ma] but we do not need them
since we have a concrete example.) The picture
below exhibits the composition

\[
	\begin{picture}(0,0)%
\epsfig{file=ps/pic2.ps}%
\end{picture}%
\setlength{\unitlength}{0.00041700in}%
%
\begingroup\makeatletter\ifx\SetFigFont\undefined
% extract first six characters in \fmtname
\def\x#1#2#3#4#5#6#7\relax{\def\x{#1#2#3#4#5#6}}%
\expandafter\x\fmtname xxxxxx\relax \def\y{splain}%
\ifx\x\y   % LaTeX or SliTeX?
\gdef\SetFigFont#1#2#3{%
  \ifnum #1<17\tiny\else \ifnum #1<20\small\else
  \ifnum #1<24\normalsize\else \ifnum #1<29\large\else
  \ifnum #1<34\Large\else \ifnum #1<41\LARGE\else
     \huge\fi\fi\fi\fi\fi\fi
  \csname #3\endcsname}%
\else
\gdef\SetFigFont#1#2#3{\begingroup
  \count@#1\relax \ifnum 25<\count@\count@25\fi
  \def\x{\endgroup\@setsize\SetFigFont{#2pt}}%
  \expandafter\x
    \csname \romannumeral\the\count@ pt\expandafter\endcsname
    \csname @\romannumeral\the\count@ pt\endcsname
  \csname #3\endcsname}%
\fi
\fi\endgroup
\begin{picture}(8670,7316)(1126,-7810)
\put(1126,-1975){\makebox(0,0)[lb]{\smash{\SetFigFont{12}{14.4}{rm}$T=$}}}
\put(6046,-2005){\makebox(0,0)[lb]{\smash{\SetFigFont{12}{14.4}{rm}$S=$}}}
\put(1951,-5365){\makebox(0,0)[lb]{\smash{\SetFigFont{12}{14.4}{rm}$T\circ_2S$}}}
\put(2731,-2231){\makebox(0,0)[lb]{\smash{\SetFigFont{12}{14.4}{rm}$*$}}}
\put(4021,-2741){\makebox(0,0)[lb]{\smash{\SetFigFont{12}{14.4}{rm}$*$}}}
\put(7936,-806){\makebox(0,0)[lb]{\smash{\SetFigFont{12}{14.4}{rm}$*$}}}
\put(3376,-746){\makebox(0,0)[lb]{\smash{\SetFigFont{12}{14.4}{rm}$*$}}}
\put(5971,-7121){\makebox(0,0)[lb]{\smash{\SetFigFont{12}{14.4}{rm}$*$}}}
\put(5341,-3986){\makebox(0,0)[lb]{\smash{\SetFigFont{12}{14.4}{rm}$*$}}}
\put(4666,-5066){\makebox(0,0)[lb]{\smash{\SetFigFont{12}{14.4}{rm}$D_1$}}}
\put(3121,-1166){\makebox(0,0)[lb]{\smash{\SetFigFont{12}{14.4}{rm}$*$}}}
\put(4996,-6671){\makebox(0,0)[lb]{\smash{\SetFigFont{12}{14.4}{rm}$*$}}}
\put(4771,-7136){\makebox(0,0)[lb]{\smash{\SetFigFont{12}{14.4}{rm}$D_5$}}}
\put(4381,-6326){\makebox(0,0)[lb]{\smash{\SetFigFont{12}{14.4}{rm}$D_4$}}}
\put(4186,-5996){\makebox(0,0)[lb]{\smash{\SetFigFont{12}{14.4}{rm}$*$}}}
\put(5911,-6701){\makebox(0,0)[lb]{\smash{\SetFigFont{12}{14.4}{rm}$D_3$}}}
\put(4921,-4586){\makebox(0,0)[lb]{\smash{\SetFigFont{12}{14.4}{rm}$*$}}}
\put(8326,-2081){\makebox(0,0)[lb]{\smash{\SetFigFont{12}{14.4}{rm}$*$}}}
\put(7846,-1256){\makebox(0,0)[lb]{\smash{\SetFigFont{12}{14.4}{rm}$*$}}}
\put(8041,-1616){\makebox(0,0)[lb]{\smash{\SetFigFont{12}{14.4}{rm}$D_4$}}}
\put(8311,-2516){\makebox(0,0)[lb]{\smash{\SetFigFont{12}{14.4}{rm}$D_5$}}}
\put(3901,-2291){\makebox(0,0)[lb]{\smash{\SetFigFont{12}{14.4}{rm}$D_3$}}}
\put(2791,-1601){\makebox(0,0)[lb]{\smash{\SetFigFont{12}{14.4}{rm}$D_1$}}}
\put(2926,-2486){\makebox(0,0)[lb]{\smash{\SetFigFont{12}{14.4}{rm}$D_2$}}}
\end{picture}

\]

The most general notion of a planar algebra that
we will contemplate is that of an algebra over
$\Bbb P$ in the sense of [Ma].  That is to say, first
of all, a disjoint union $V_k$ of vector spaces
for $k>0$ and two vector spaces
$V_0^{\text{white}}$ and $V_0^{\text{black}}$
(which we will call $P_0$ and $P_{1,1}$ later on).
Linear maps between tensor powers of these vector
spaces form a coloured operad \ Hom \ in the
obvious way under composition of maps and the
planar algebra structure  on the $V$'s is given by
a morphism of coloured operads from $\Bbb P$ to Hom. 
In
practice this means that, to a $k$-tangle $T$ in
$\Bbb P$ there is a linear map
$Z(T):\bigotimes^n_{i=1}V_{k_i}\to V_k$ such that
$Z(T\circ_i S)=Z(T)\circ_i Z(S)$ where the
$\circ_i$ on the right-hand side is composition of
linear maps in \ Hom \ .

Note that the vector spaces $V_0^{\text{white}}$ and $V_0^{\text{black}}$
may be different. This is the case for the ``spin
models" of $\S$3. Both these $V_0$'s become
commutative associative algebras using the tangles

\[
	\begin{picture}(0,0)%
\epsfig{file=ps/pic3.ps}%
\end{picture}%
\setlength{\unitlength}{0.00041700in}%
%
\begingroup\makeatletter\ifx\SetFigFont\undefined
% extract first six characters in \fmtname
\def\x#1#2#3#4#5#6#7\relax{\def\x{#1#2#3#4#5#6}}%
\expandafter\x\fmtname xxxxxx\relax \def\y{splain}%
\ifx\x\y   % LaTeX or SliTeX?
\gdef\SetFigFont#1#2#3{%
  \ifnum #1<17\tiny\else \ifnum #1<20\small\else
  \ifnum #1<24\normalsize\else \ifnum #1<29\large\else
  \ifnum #1<34\Large\else \ifnum #1<41\LARGE\else
     \huge\fi\fi\fi\fi\fi\fi
  \csname #3\endcsname}%
\else
\gdef\SetFigFont#1#2#3{\begingroup
  \count@#1\relax \ifnum 25<\count@\count@25\fi
  \def\x{\endgroup\@setsize\SetFigFont{#2pt}}%
  \expandafter\x
    \csname \romannumeral\the\count@ pt\expandafter\endcsname
    \csname @\romannumeral\the\count@ pt\endcsname
  \csname #3\endcsname}%
\fi
\fi\endgroup
\begin{picture}(7545,3028)(871,-4570)
\put(2161,-3821){\makebox(0,0)[lb]{\smash{\SetFigFont{12}{14.4}{rm}$D_2$}}}
\put(2146,-2621){\makebox(0,0)[lb]{\smash{\SetFigFont{12}{14.4}{rm}$D_1$}}}
\put(4396,-3161){\makebox(0,0)[lb]{\smash{\SetFigFont{12}{14.4}{rm}and}}}
\put(6676,-3821){\makebox(0,0)[lb]{\smash{\SetFigFont{12}{14.4}{rm}$D_2$}}}
\put(6661,-2621){\makebox(0,0)[lb]{\smash{\SetFigFont{12}{14.4}{rm}$D_1$}}}
\end{picture}

\]

To handle tangles with no internal discs we decree
that the tensor product over the empty set be the
field $K$ and identify Hom$(K,V_k)$ with $V_k$ so
that each $V_k$ will contain a privileged subset
which is $Z(\{ k$-tangles with no internal
discs$\})$. This is the ``unital"
 structure (see [Ma]).

One may want to impose various conditions such as
dim$(V_k)<\i$ for all $k$. The condition
dim$(V_0^{\text{white}})=1=$dim$(V_0^{\text{black}})$
is significant and we impose it in our formal
definition of planar algebra (as opposed to
general planar algebra) later on.  It implies that
there is a unique way to identify each $V_0$ with
$K$ as algebras, and $Z(\bigcirc)=1=Z(\deyea )$.
There are thus also two scalars associated to a
planar algebra, $\d_1=Z(\weye )$ and
$\d_2=Z(\beye )$ (the inner circles are strings,
not discs!). It follows that  $Z$ is
multiplicative on connected components, i.e., if a
part of a tangle $T$ can be surrounded by a disc
so that $T=T'\circ_i S$ for a tangle $T'$ and a
0-tangle $S$, then $Z(T)=Z(S)Z(T')$ where $Z(S)$
is a multilinear map into the field $K$.

Two simple examples serve as the keys to
understanding the notion of a planar algebra. The
first is the Temperley-Lieb algebra $TL$, some
vestige of which is present in every planar
algebra. The vector spaces $TL_k$ are:
$$
TL_0^{\text{black}}\simeq
TL_0^{\text{white}}\simeq K
$$
and $TL_k$ is the vector space  whose basis is the
set of diffeomorphism classes of connected planar
$k$-tangles with no internal discs, of which there
are $\frac{1}{k+1}{2k\choose k}$. The
action of a planar $k$-tangle on $TL$ is almost
obvious --- when one fills the internal discs of a
tangle with basis elements of $TL$ one obtains
another basis element, except for some simple
closed curves. Each closed curve counts a
multiplicative factor of $\d$ and then is removed.
It is easily verified that this defines an action
of $\Bbb P$ on $TL$.
As we have observed, any planar algebra contains
elements corresponding to the $TL$ basis. They are
not necessarily linearly independent. See [GHJ]
and $\S$2.1.

 The second key example of a planar algebra is
given by tensors. We think of a tensor as an
object which yields a number each time its indices
are specified. Let $V_k$ be the vector space of
tensors with $2k$ indices. An element of $\Bbb P$
gives a scheme for contracting tensors, once a
tensor is assigned to each internal disc. The
indices lie on the strings and are locally
constant thereon. The boundary indices are fixed
and are the indices of the output tensor. All
indices  on strings not touching $D$ are summed
over and one contracts by taking, for a given set
of indices, the product of the values of the
tensors in the internal discs. One recognizes the
partition function of a statistical mechanical
model ([Ba]), the boundary index values being the
boundary conditions and the tensor values being
the Boltzmann weights. This diagrammatic
contraction calculus for tensors is well known
([Pe]) but here we are only considering {\it
planar} contraction systems. If the whole planar
algebra of all tensors were the only example this
subject would be of no interest, but in fact there
is a huge family of planar subalgebras --- vector
spaces of tensors closed under planar contractions
--- rich enough to contain the theory of finitely
generated groups and their Cayley graphs. See
$\S$2.7.

The definition of planar algebra we give in $\S$1
is not the operadic one. When the planar algebra
structure first revealed itself, the $V_k$'s
already had an associative algebra structure
coming from the von Neumann algebra context. Thus
our definition will be in terms of a universal planar
algebra on some set of generators (labels) which
can be combined in arbitrary planar fashion. The
discs we have used above become boxes in section one, reflecting
the  specific algebra structure we began with. The
equivalence of the two definitions is completed in
Proposition 1.20. The main ingredient of the equivalence is that the
planar operad $\Bbb P$ is generated by the Temperley-Lieb algebra
and tangles of two kinds:
\begin{enumerate}
\item Multiplication, which is the following
tangle (illustrated for k =5)

\[
	\begin{picture}(0,0)%
\epsfig{file=ps/pic4.ps}%
\end{picture}%
\setlength{\unitlength}{0.00041700in}%
%
\begingroup\makeatletter\ifx\SetFigFont\undefined
% extract first six characters in \fmtname
\def\x#1#2#3#4#5#6#7\relax{\def\x{#1#2#3#4#5#6}}%
\expandafter\x\fmtname xxxxxx\relax \def\y{splain}%
\ifx\x\y   % LaTeX or SliTeX?
\gdef\SetFigFont#1#2#3{%
  \ifnum #1<17\tiny\else \ifnum #1<20\small\else
  \ifnum #1<24\normalsize\else \ifnum #1<29\large\else
  \ifnum #1<34\Large\else \ifnum #1<41\LARGE\else
     \huge\fi\fi\fi\fi\fi\fi
  \csname #3\endcsname}%
\else
\gdef\SetFigFont#1#2#3{\begingroup
  \count@#1\relax \ifnum 25<\count@\count@25\fi
  \def\x{\endgroup\@setsize\SetFigFont{#2pt}}%
  \expandafter\x
    \csname \romannumeral\the\count@ pt\expandafter\endcsname
    \csname @\romannumeral\the\count@ pt\endcsname
  \csname #3\endcsname}%
\fi
\fi\endgroup
\begin{picture}(3759,3756)(661,-4406)
\put(2319,-1916){\makebox(0,0)[lb]{\smash{\SetFigFont{12}{14.4}{rm}$D_1$}}}
\put(2318,-3417){\makebox(0,0)[lb]{\smash{\SetFigFont{12}{14.4}{rm}$D_2$}}}
\put(1718,-1039){\makebox(0,0)[lb]{\smash{\SetFigFont{12}{14.4}{rm}$*$}}}
\put(1823,-1624){\makebox(0,0)[lb]{\smash{\SetFigFont{12}{14.4}{rm}$*$}}}
\put(1823,-3048){\makebox(0,0)[lb]{\smash{\SetFigFont{12}{14.4}{rm}$*$}}}
\end{picture}

\]

\item Annular tangles: ones with only one internal
disc, e.g.,
\end{enumerate}

\[
	\begin{picture}(0,0)%
\epsfig{file=ps/pic5.ps}%
\end{picture}%
\setlength{\unitlength}{0.00041700in}%
%
\begingroup\makeatletter\ifx\SetFigFont\undefined
% extract first six characters in \fmtname
\def\x#1#2#3#4#5#6#7\relax{\def\x{#1#2#3#4#5#6}}%
\expandafter\x\fmtname xxxxxx\relax \def\y{splain}%
\ifx\x\y   % LaTeX or SliTeX?
\gdef\SetFigFont#1#2#3{%
  \ifnum #1<17\tiny\else \ifnum #1<20\small\else
  \ifnum #1<24\normalsize\else \ifnum #1<29\large\else
  \ifnum #1<34\Large\else \ifnum #1<41\LARGE\else
     \huge\fi\fi\fi\fi\fi\fi
  \csname #3\endcsname}%
\else
\gdef\SetFigFont#1#2#3{\begingroup
  \count@#1\relax \ifnum 25<\count@\count@25\fi
  \def\x{\endgroup\@setsize\SetFigFont{#2pt}}%
  \expandafter\x
    \csname \romannumeral\the\count@ pt\expandafter\endcsname
    \csname @\romannumeral\the\count@ pt\endcsname
  \csname #3\endcsname}%
\fi
\fi\endgroup
\begin{picture}(3000,3018)(308,-3695)
\put(1193,-1076){\makebox(0,0)[lb]{\smash{\SetFigFont{12}{14.4}{rm}$*$}}}
\put(1351,-3061){\makebox(0,0)[lb]{\smash{\SetFigFont{12}{14.4}{rm}$*$}}}
\end{picture}

\]

The universal planar algebra is useful for
constructing planar algebras and restriction to
the two generating tangles sometimes makes it
shorter to check that a given structure is a
planar  algebra.
The algebra structure we begin with in $\S$1 corresponds of
course to the multiplication tangle given above.

The original algebra structure has been studied in
some detail (see $\S$3.1) but it should be quite
clear that the operad provides a vast family of
algebra structures on a planar algebra which we
have only just begun to appreciate. For instance,
the annular tangles above form an algebra over
which all the $V_k$'s in a planar algebra are
modules. This structure alone seems quite rich
([GL]) and we exploit it just a little to get
information on principal graphs of subfactors in
4.2.11. We have obtained more sophisticated
results in terms of generating functions  which we
will present in a future paper.

We present several examples of planar algebras in
$\S$2, but it is the connection with subfactors
that has been our main motivation and guide for
this work. The two leading theorems occur in
$\S$4. The first one shows how to obtain a planar
algebra from  a finite index subfactor $N\subset
M$. The vector space $V_k$ is the set of
$N$-central vectors in the $N-N$ bimodule
$M_{k-1}=M\otimes_N M\otimes_N\dots\otimes_N M$
($k$ copies of $M$), which, unlike $M_{k-1}$
itself, is finite dimensional. The planar algebra
structure on these $V_k$'s is obtained by a method
reminiscent of topological quantum field theory.
Given a planar  $k$-tangle $T$ whose internal
discs are labelled by elements of the $V_j$'s, we
have to show how to construct an element of $V_k$,
associated with the boundary of $T$, in a natural
way. One starts with a very small circle (the
``bubble") in the distinguished white region of
$T$, tangent to the boundary of $D$. We then allow
this circle to bubble out until it gets to the
boundary. On its way the bubble will have to cross
strings of the tangle and envelop internal discs.
As it does so it acquires shaded intervals which
are its intersections  with the shaded regions of
$T$. Each time the bubble envelops an internal
disc $D_i$, it acquires $k_i$ such shaded
intervals and, since an element of $V_{k_i}$ is a
tensor in $\otimes^{k_i}_N M$, we assign elements
of $M$ to the shaded region according to this
tensor.  There are also rules for assigning and
contracting tensors as the bubble crosses strings
of the tangle. At the end we have an element of
$\otimes^k_N M$ assigned to the boundary. This is
the action of the operad element on the vectors in
$V_{k_i}$. Once the element of $\otimes^k_N M$ has
been constructed and shown to be invariant under
diffeomorphisms, the formal operadic properties
are immediate.

One could try to carry out this procedure for an
arbitrary inclusion $A\subset B$ of rings, but
there are a few obstructions involved in showing
that our bubbling process is well defined. Finite
index (extremal) subfactors have all the special
properties required, though there are surely other
families of subrings for which the procedure is
possible.

The following tangle:

\[
	\begin{picture}(0,0)%
\epsfig{file=ps/pic6.ps}%
\end{picture}%
\setlength{\unitlength}{0.00041700in}%
%
\begingroup\makeatletter\ifx\SetFigFont\undefined
% extract first six characters in \fmtname
\def\x#1#2#3#4#5#6#7\relax{\def\x{#1#2#3#4#5#6}}%
\expandafter\x\fmtname xxxxxx\relax \def\y{splain}%
\ifx\x\y   % LaTeX or SliTeX?
\gdef\SetFigFont#1#2#3{%
  \ifnum #1<17\tiny\else \ifnum #1<20\small\else
  \ifnum #1<24\normalsize\else \ifnum #1<29\large\else
  \ifnum #1<34\Large\else \ifnum #1<41\LARGE\else
     \huge\fi\fi\fi\fi\fi\fi
  \csname #3\endcsname}%
\else
\gdef\SetFigFont#1#2#3{\begingroup
  \count@#1\relax \ifnum 25<\count@\count@25\fi
  \def\x{\endgroup\@setsize\SetFigFont{#2pt}}%
  \expandafter\x
    \csname \romannumeral\the\count@ pt\expandafter\endcsname
    \csname @\romannumeral\the\count@ pt\endcsname
  \csname #3\endcsname}%
\fi
\fi\endgroup
\begin{picture}(3004,3012)(304,-3679)
\put(1846,-1609){\makebox(0,0)[lb]{\smash{\SetFigFont{12}{14.4}{rm}$*$}}}
\put(1441,-919){\makebox(0,0)[lb]{\smash{\SetFigFont{12}{14.4}{rm}$*$}}}
\end{picture}

\]

\noindent defines a rotation of period $k$ ($2k$
boundary points) so it is a consequence of the
planar algebra structure that the rotation
$x_1\otimes x_2\otimes\dots\otimes x_k\mapsto
x_2\otimes x_3\otimes\dots\otimes x_k\otimes x_1$,
which makes no sense on $M_{k-1}$, is well defined
on $N$-central vectors and has period $k$. This
result is in fact an essential technical
ingredient of the proof of Theorem 4.2.1.

Note that we seem to have avoided the use of
correspondences in the sense of Connes ([Co1]) by
working in the {\it purely algebraic} tensor
product. But the avoidance of $L^2$-analysis,
though extremely convenient, is a little illusory
since the proof of the existence and periodicity
of the rotation uses $L^2$ methods. The Ocneanu
approach ([EK]) uses the $L^2$ definition and the vector spaces $V_k$ are defined
as $Hom_{N,N}(\otimes^j_N M)$ and $Hom_{N,M}(\otimes^j_N M)$
depending on the parity of k. It is no doubt possible
to give a direct proof of Theorem 4.2.1 using this
definition - this would be the ``hom" 
version, our method being the ``$\otimes$" 
method.

To identify the operad structure with the usual
algebra structure on $M_{k-1}$ coming from the
``basic construction" of [J1], we show that the
multiplication tangle above does indeed define the
right formula. This, and a few similar details, is
suprisingly involved and accounts for some
unpleasant looking formulae in $\S$4. Several other
subfactor notions, e.g.~tensor product, are shown
to correspond to their planar algebra
counterparts, already abstractly defined in $\S$3.
Planar algebras also inspired, in joint work with
D.~Bisch, a notion of free product. We give the
definition here and will explore this notion in a
forthcoming paper with Bisch.

The second theorem of $\S$4 shows that one can
construct a subfactor from a planar algebra with
$*$-structure and suitable ``reflection"
positivity. It is truly remarkable that the axioms
needed by Popa for his construction of subfactors
in [Po2] follow so closely the axioms of planar
algebra, at least as formulated using boxes and
the universal planar algebra. For Popa's construction is
quite different from the ``usual" one of
[J1], [F+], [We1], [We2]. Popa uses an
amalgamated free product construction which
introduces an unsatisfactory element in the
correspondence between planar algebras and
subfactors. For although it is true that the
standard invariant of Popa's subfactor  is indeed
the planar algebra from which the subfactor was
constructed, it is {\it not} true that, if one
begins with a subfactor $N\subset M$, even
hyperfinite, and applies Popa's procedure to the
standard invariant, one obtains $N\subset M$ as a
result. There are many difficult questions here,
the main one of which is to decide when a given
planar algebra arises from a subfactor of the
Murray-von Neumann hyperfinite type II${}_1$
factor ([MvN]).

There is a criticism that has and should be made of 
our definition of a planar algebra - that it is
too restrictive. By enlarging the class of tangles in the   
planar operad, say so as to include oriented edges and 
boundary points, or discs with an odd number of 
boundary points, one would obtain a notion of planar
algebra applicable to more examples. For instance,
if the context were the study of group representations
our definition would have us studying say, SU(n) by 
looking at tensor powers of the form ${V\otimes\bar V\otimes V\otimes \bar V\dots}$
(where V is the defining representation on $\Bbb C^n$)



